\section{Introduction}\label{sec:1}

\subsection{Central cuts, ellipsoid computation, and rounding}\label{sec:1.1}

A convex body is a compact convex subset of $\mathbb R^n$ with nonempty interior. Let $J(K)$ be its unique maximum-volume inscribed ellipsoid, let $c(K)$ be its center, and set
\begin{equation}\label{eq:1.1}
w(K)=\operatorname{vol}J(K)
=\max\{\operatorname{vol}E:E\subseteq K,\ E\text{ is an ellipsoid}\},
\end{equation}
where $\operatorname{vol}$ denotes Lebesgue volume. Existence, uniqueness, and the contact characterization are classical \cite{John1948,GG2007}. We study the change in $w$ under a halfspace cut whose boundary contains $c(K)$.

Such cuts arise in localization methods for convex optimization. An oracle supplies a halfspace containing the target set, and the current region is replaced by its intersection with that halfspace. The choice of center controls the geometric progress. In the centroid method, Grünbaum's theorem bounds the surviving ordinary-volume fraction by $1-(n/(n+1))^n<1-e^{-1}$ \cite{Grunbaum1960}, but computing the centroid is expensive. The classical enclosing-ellipsoid method admits an $O(n^2)$ arithmetic update, with a dimension-dependent contraction factor \cite{GLS1988,Todd2023}. Analytic and volumetric centers offer other computational tradeoffs \cite{BV2004,AV1995,Vaidya1996,Anstreicher1997}.

The method of inscribed ellipsoids instead uses $c(K)$ and the potential $w(K)$ \cite{TKE1988}. It also provides a geometric normalization. If $J(K)=c+AB_2^n$, John's rounding theorem gives
\begin{equation}\label{eq:1.2}
c+AB_2^n\subseteq K\subseteq c+nAB_2^n,
\qquad
B_2^n\subseteq A^{-1}(K-c)\subseteq nB_2^n.
\end{equation}

For a centrally symmetric body, the factor $n$ improves to $\sqrt n$. This rounding is useful in geometric preconditioning, volume approximation, and fixed-dimensional integer programming \cite{GG2007,KT1993}.

For a polytope with $N$ inequalities, computing $J(K)$ is a finite convex optimization problem; Section~\ref{sec:2} gives its formulation and dual certificates. Interior-point methods exploit its structure \cite{NN1994,KT1993,ZG2003}. One established bound computes an ellipsoid of volume at least $e^{-\delta}w(K)$, $0<\delta<1$, under a known normalization $B_2^n\subseteq K\subseteq\Lambda B_2^n$, within $O(N^{7/2}\log(N\Lambda/\delta))$ arithmetic operations \cite{Anstreicher2002}. This is a bound for approximate ellipsoid computation. The contraction theorem below concerns the exact John center; relating a computed center to an admissible cut requires the corresponding approximation analysis. Section~\ref{sec:7} distinguishes this computational cost from the number of exact central cuts.

Khachiyan asked for the optimal dimension-independent contraction factor \cite{Khachiyan1990}. Zhou, Zou, and Liu proved the value $\sqrt e/2$ and its uniform optimality \cite{Zhou2026}. The present paper determines the exact value in each dimension and describes the geometry of equality and near equality.

\subsection{Contributions}\label{sec:1.2}

The results follow two connected threads.

\begin{enumerate}[label=\arabic*.,leftmargin=*]

\item \textbf{Sharp central-cut contraction.} Let $C_n$ be the supremum of $w(K\cap H)/w(K)$ over convex bodies and halfspaces through their John centers. We prove
\begin{equation}\label{eq:1.3}
C_n=R_n,\qquad R_1=\frac12,\qquad
R_n=\frac12\left(\frac n{n-1}\right)^{(n-1)/2}\quad(n\ge2),
\end{equation}
with attainment in every dimension. The same sharp constant is obtained by restricting the supremum to finite John contact polytopes: the reduction in Section~\ref{sec:2.4} preserves the original ellipsoid volume while enlarging the cut region. The proof constructs an explicit conic certificate from the resulting finite data and bounds its determinant. A simple extremal configuration determines the candidate ellipsoid and the certificate parameters.

\item \textbf{Equality and stability from an exact deficit.} The logarithmic shortfall from $R_n$ is the sum of a geometric determinant deficit, a cone-complementarity term, and a positive-definite matrix divergence. Their zero sets give all equality cases, including nonsimplex extremizers. Their quantitative control yields dimension-independent bounds for the new ellipsoid's shape and center and for the weighted contact geometry.

\end{enumerate}

John's theorem and the finite convex optimization model are background. The problem-specific contribution is the quantitative certificate estimate and the geometric structure that it reveals.

\subsection{Proof organization}\label{sec:1.3}

Section~\ref{sec:2} develops the optimization framework, proves exactness of the
certificate, and reduces the geometry to finite John contact data. Section~\ref{sec:3}
proves the sharp contraction using the frame estimate. Sections~\ref{sec:4}--\ref{sec:6} derive
the exact deficit, equality, and stability; Section~\ref{sec:7} gives iteration and
computational consequences. Section~\ref{sec:8} proves the determinant estimate,
with higher dimensions before the low-dimensional block analysis, and
Section~\ref{sec:9} supplies the stability estimates. Section~\ref{sec:10} discusses extensions
and further matrix properties. Section~\ref{sec:11} concludes.
